% GENERATED FILE. Edit canonical lessons, then run npm run generate:books.
% canonical-source-hash: fnv1a64:1dad939897dbe049
% canonical-lessons: MW-C159-gharai, MW-C159-gayo, MW-R159-went-home

\chapter{Went Home}
\label{ch:mw-a2-went-home}
\hlchaptermodality{159}

\begin{chapteropening}
\textbf{By the end of this chapter:} \emph{I can say that someone is at home, and that I or someone else went somewhere.}

Complete MW-R159-went-home: say who is at home and who went where, with the ending that agrees.
\end{chapteropening}

\section[\texorpdfstring{\mw{घरै} (gharai)}{घरै (gharai)}]{\mw{घरै} — at home}

\label{lesson:MW-C159-gharai}



\begin{quote}
Say \emph{here}, \emph{there} and \emph{on top}: \textbf{\mw{अठै}}, \textbf{\mw{बठै}}, \textbf{\mw{माथै}}. Then \emph{a house}: \textbf{\mw{घर}}.
\end{quote}

\begin{grammarlens}[title={Grammar Lens: \mw{घर} with -\mw{ै}}]
\textbf{\mw{घर}} with the \textbf{-\mw{ै}} of \textbf{\mw{अठै}}, \textbf{\mw{बठै}} and \textbf{\mw{माथै}} is \textbf{\mw{घरै}}: at home, or home as the place you go to.

\begin{quote}
\textbf{\mw{हूं} \mw{घरै} \mw{हूं।}}
\end{quote}

\begin{quote}
\emph{hū̃ gharai hū̃.} — I am at home.
\end{quote}

\begin{quote}
\textbf{\mw{वो} \mw{रोज} \mw{आथण} \mw{घरै} \mw{आवै} \mw{है।}}
\end{quote}

\begin{quote}
\emph{vo roj āthaṇ gharai āvai hai.} — He comes home every evening.
\end{quote}

Old Marwari had the same ending in \emph{āgai}, \textquotedblleft{}before\textquotedblright{} (CDIAL 68).
\end{grammarlens}

\subsection*{Your turn}
Say these aloud:

\begin{itemize}
\raggedright
  \item \textbf{\mw{मां} \mw{घरै} \mw{है।}} — \emph{mā̃ gharai hai.} — Mother is at home.
  \item \textbf{\mw{आप} \mw{रोज} \mw{घरै} \mw{जावो} \mw{हो}?} — \emph{āp roj gharai jāvo ho?} — Do you go home every day?
  \item \textbf{\mw{हां} \mw{सा}, \mw{हूं} \mw{रोज} \mw{घरै} \mw{जाऊं} \mw{हूं।}} — \emph{hā̃ sā, hū̃ roj gharai jāū̃ hū̃.} — Yes, I go home every day.
\end{itemize}

\begin{tcolorbox}[breakable,colback=teal!4,colframe=teal!35!black,title={Before you move on}]
\emph{At home}? (\textbf{\mw{घरै}}.) \emph{Here}? (\textbf{\mw{अठै}}.)
\end{tcolorbox}
\section[\texorpdfstring{\mw{गयो} · \mw{गई} · \mw{गया} (gayo · gaī · gayā)}{गयो · गई · गया (gayo · gaī · gayā)}]{\mw{गयो}, \mw{गई}, \mw{गया} — went}

\label{lesson:MW-C159-gayo}



\begin{quote}
Say \emph{I go to the market every day}. Then \emph{yesterday}, \textbf{\mw{काल}}, and \emph{a boy}, \emph{a girl}.
\end{quote}

\begin{grammarlens}[title={Grammar Lens: one word for the past}]
\textbf{\mw{गयो}} is the past of \textbf{\mw{जाणो}}, all on its own: no \textbf{\mw{हूं}} and no \textbf{\mw{है}} after it.

\begin{quote}
\textbf{\mw{हूं} \mw{काल} \mw{बाजार} \mw{गयो।}}
\end{quote}

\begin{quote}
\emph{hū̃ kāl bājār gayo.} — I went to the market yesterday.
\end{quote}

It ends like \textbf{\mw{छोरो}}, and changes like it, to agree with whoever went: \textbf{\mw{गयो}} for a man, \textbf{\mw{गई}} for a woman, \textbf{\mw{गया}} for more than one.

\begin{quote}
\textbf{\mw{मां} \mw{घरै} \mw{गई।}}
\end{quote}

\begin{quote}
\emph{mā̃ gharai gaī.} — Mother went home.
\end{quote}

\begin{quote}
\textbf{\mw{सगळा} \mw{घरै} \mw{गया।}}
\end{quote}

\begin{quote}
\emph{sagḷā gharai gayā.} — Everyone went home.
\end{quote}

With a past, \textbf{\mw{काल}} is \emph{yesterday}. CDIAL lists Old Marwari \emph{gayaü}, \textquotedblleft{}gone\textquotedblright{} (CDIAL 4008).
\end{grammarlens}

\subsection*{Your turn}
A woman says \textbf{\mw{हूं} \mw{काल} \mw{बाजार} \mw{गई।}} (\emph{hū̃ kāl bājār gaī.})

Say these aloud:

\begin{itemize}
\raggedright
  \item \textbf{\mw{राम} \mw{काल} \mw{बाजार} \mw{गयो।}} — \emph{rām kāl bājār gayo.} — Ram went to the market yesterday.
  \item \textbf{\mw{मां} \mw{घरै} \mw{गई।}} — \emph{mā̃ gharai gaī.} — Mother went home.
  \item \textbf{\mw{वो} \mw{कोनी} \mw{गयो।}} — \emph{vo konī gayo.} — He didn't go.
\end{itemize}

\begin{tcolorbox}[breakable,colback=teal!4,colframe=teal!35!black,title={Before you move on}]
\emph{He went}? (\textbf{\mw{गयो}}.) \emph{She went}, of mother? (\textbf{\mw{गई}}.) \emph{Everyone went}? (\textbf{\mw{गया}}.)
\end{tcolorbox}
\section[(dialogue)]{Went home — a review}

\label{lesson:MW-R159-went-home}



\begin{quote}
Say \emph{I am at home}.
\end{quote}

\subsection*{Your turn}
\noindent
\begin{tabularx}{\linewidth}{@{}>{\raggedright\arraybackslash}X>{\raggedright\arraybackslash}X@{}}
\toprule
\textbf{who went} & \textbf{went} \\
\midrule
\textbf{\mw{राम}}, a man & \textbf{\mw{गयो}} \\
\textbf{\mw{मां}}, a woman & \textbf{\mw{गई}} \\
\textbf{\mw{सगळा}}, more than one & \textbf{\mw{गया}} \\
\bottomrule
\end{tabularx}

Say these aloud:

\begin{itemize}
\raggedright
  \item \emph{I went to the market yesterday}, as a man and then as a woman
  \item \emph{mother went home}; \emph{everyone went home}
  \item \emph{he didn't go}
  \item \emph{do you go home every day?}; \emph{yes, I go home every day}
\end{itemize}

\begin{tcolorbox}[breakable,colback=teal!4,colframe=teal!35!black,title={Before you move on}]
\emph{At home}? (\textbf{\mw{घरै}}.) \emph{Went}, for a man, a woman, and more than one? (\textbf{\mw{गयो}}, \textbf{\mw{गई}}, \textbf{\mw{गया}}.)
\end{tcolorbox}
